% Draft: research question (with literature motivation) + answer. REFRAMED.
% For review -- NOT inserted into the manuscript.
%
% Reframing (v2): the identification claim is NOT an aggregation argument. The
% classification of which statistics reveal the sequence effect is settled by
% the two reading sequences themselves, before any averaging (manuscript's
% "settled before averaging", Prop PRO), and is benchmark-free (Prop ORD). So
% this version centres the between-sequence contrast on a single statistic and
% explicitly DEMOTES both aggregation and the benchmark to non-load-bearing.
%
% Terminology is the manuscript's: sequence / sequence effect (def:seqeffect),
% protected/unprotected (def:protection), cross-attribute association
% (def:assoc), first/second order in c, benchmark PB. Proposed additions:
% "sequence effect on a statistic" (F(PJ_AB)-F(PJ_BA)) and Prop ORD.

% ---------- QUESTION (with literature motivation) ----------

That coherent updating on two uncertain impressions depends on the sequence in
which they arrive is not in question. Jeffrey conditioning is sequence-dependent,
and the condition under which the two sequences agree---that neither cue disturb
the other's partition---has been known since \citet{DiaconisZabell1982} and was
recast by \citet{HawthorneAND2004} as the commutation condition for standard
sequential updating. That literature settled \emph{whether} the two sequences
coincide, reasoning throughout about the evaluator's whole joint belief. It
leaves untouched a prior question about measurement: when the two sequences
differ, does the difference survive into the quantities one would actually read
off an evaluator---a level of belief in one attribute, a decision, the believed
cross-attribute association (Definition~\ref{def:assoc})---or does it vanish from
some of them? The sequence effect (Definition~\ref{def:seqeffect}) is a
difference between two whole joint beliefs; whether it persists into a given
summary of those beliefs is a separate matter, and it is the matter that decides
whether the effect can be seen at all.

This is not a question about pooling many evaluators together. It is a question
about the summary itself. Fix any statistic $F$ of the belief and compare its
value for an evaluator who read the cues in one sequence with its value for one
who read them in the other. For some statistics this difference is of the same
order in the prior covariance $c$ as the sequence effect itself---large enough
to clear any fixed measurement precision---while for others it is of strictly
higher order, below resolution however precisely one measures
(Definition~\ref{def:protection}). This paper asks which statistics fall on
which side of that line and, because the line turns out to be sharp, what an
observer who reads only a statistic on the wrong side of it can and cannot infer.

% ---------- ANSWER (triangular: plain finding first, mechanism after) ----------

The answer is that the sequence survives into some statistics and vanishes from
others, and which is which is fully determined. The believed cross-attribute
association---the statistic most naturally read as the signature of a
stereotype---does not carry it: two evaluators who met the same impressions in
opposite sequences report essentially the same association, differing only at
second order in $c$. The marginal probabilities and the decision do carry it,
differing already at first order. So one evaluator's beliefs can show a plain
sequence effect in their marginals and their decision while their believed
association looks sequence-free.

This is cautionary for measurement, and it is exactly the loop the question left
open. A test that asks only whether an evaluator's believed association matches
its sequence-free value will report conformity even when the sequence effect is
plainly present in the same beliefs; a null on the association is then evidence
about the choice of instrument, not about the evaluator. The statistics that do
carry the effect are the marginals and the decision. This sharpens what the
documented order effects of impression formation
\citep{HogarthEinhorn1992,Asch1946} can establish: those studies elicit a single
evaluative level, which is a carrying statistic and will show the effect, but a
level on its own cannot separate a sequence-dependent evaluator from a
sequence-free one---that separation needs the association read alongside it. The
instrument the effect calls for is a carrying statistic---a marginal, or the
decision---compared between the two sequences, with the association elicited
alongside it. The same comparison identifies the mechanism as well as the
effect: the marginal that stays fixed as the prior association varies sits on
the last-read attribute under overwriting, and on the first-read attribute
under a first-impression rule.

The classification is this sharp because it is decided by the two sequences
alone. Each sequence pushes the belief in one characteristic direction; a
statistic carries the sequence effect exactly when it varies along the span of
those two directions, and loses it when its sensitivity at the independent belief
points instead along the believed association. That condition is necessary as
well as sufficient (Proposition~\ref{prop:PRO}), so the statistics that lose the
effect are precisely the local dependence measures---the association, the odds
ratio, Yule's $Q$---together with the surplus-weighted loss for the threshold
decision (Theorem~\ref{thm:LOS}); every other smooth statistic carries it at
first order (Propositions~\ref{prop:DIV},~\ref{prop:DRF}; Lemma~\ref{lem:SCR}).

Two things this classification does not rest on. It does not rest on pooling. The
line is drawn before any averaging over a population of evaluators takes place,
so gathering individual rather than aggregate data neither reveals a lost effect
nor hides a carried one---a statistic on the wrong side of the line stays there
at every level of aggregation, which is why a poorly chosen statistic cannot be
rescued by finer data. And it does not rest on the benchmark. A statistic loses
the sequence effect exactly when the effect on it,
$F(\PJ_{AB})-F(\PJ_{BA})$, is of second order (Proposition~\ref{prop:ORD})---a
comparison of the two sequences to each other that names neither the
sequence-free benchmark $\PB$ nor any mixture of the sequences. The benchmark
re-enters only as a bridge to the prior literature: the sequence-free value the
losing statistics conform to is precisely the Bayes-factor rule of
\citet{Field1978} and \citet{Wagner2002}.
